\mychapter{II. LITERATURE REVIEW}{II. LITERATURE REVIEW}

\hs Several research fields come together in this project: natural language processing,
information retrieval, and the growing area of legal technology. This chapter reviews the
work most relevant to the system proposed here. It moves from the general idea of
retrieval-augmented generation, through the retrieval methods that make it work, toward
the more recent developments in query rewriting, legal question answering, agentic
reasoning, and the particular difficulties of low-resource languages such as Khmer.

\section{Retrieval-Augmented Generation (RAG)}

\hs The idea of Retrieval-Augmented Generation was introduced by \citet{lewis2020rag}.
Rather than relying on a model to recall facts from its parameters alone, their framework
first retrieves relevant documents from a dense vector index and then conditions
generation on those documents. On open-domain question answering benchmarks, this
combination outperformed models that depended only on parametric memory. The significance
of the result was that grounding generation in retrieved evidence became a practical and
repeatable recipe. It is the same recipe that this thesis adapts for the legal domain,
where an answer that cannot be traced back to a specific provision is of little use.

\section{Dense and Sparse Retrieval Methods}

\hs Retrieval comes in two broad forms, and both matter here. The dense form is well
represented by Dense Passage Retrieval, or DPR, from \citet{karpukhin2020dpr}. DPR
encodes queries and passages with bi-encoder neural networks into a shared vector space,
and on open-domain QA it clearly improved over keyword matching. The same work is candid
about a weakness, however. When a query hinges on a rare or very specific term, dense
retrieval can miss it, because such terms are thinly represented in the learned space.
For legal text in a low-resource language like Khmer, that weakness is hard to ignore.

\hs The sparse form has a much longer history. \citet{robertson2009bm25} formalised
BM25, which scores documents using term frequency, inverse document frequency, and a
correction for document length. It is a simple model, yet it has stayed competitive for
decades and still powers many production search systems. BM25 is a natural fit for legal
retrieval for one plain reason: legal language is exact. Whether a specific legal term
actually appears in an article often tells you more than how semantically similar the
article feels.

\hs On the engineering side, \citet{johnson2021faiss} developed FAISS, a library built
for fast similarity search across large sets of dense vectors. It supports both exact
search and approximate methods, which makes dense retrieval feasible at scale. The system
in this thesis uses FAISS with a flat inner product index, giving exact cosine similarity
over the indexed article embeddings.

\section{Hybrid Retrieval and Rank Fusion}

\hs Reading the dense and sparse literature side by side suggests an obvious move:
combine them. The two approaches fail in different ways. Dense retrieval understands
meaning even when the wording changes, as \citet{karpukhin2020dpr} showed, while BM25
catches the exact terms that dense methods sometimes overlook, as the durability of
\citet{robertson2009bm25} attests. Their weaknesses do not overlap, so together they
cover more ground than either does alone. This is especially helpful when a casual Khmer
question and a formal legal article barely share any vocabulary.

\hs The remaining question is how to merge two ranked lists, and here this work follows
\citet{cormack2009rrf}. Their Reciprocal Rank Fusion combines rankings without having to
normalise scores across systems, using the formula
$\text{RRF}(d) = \sum_{r \in R} \frac{1}{k + r(d)}$, where $k$ is a smoothing constant
and $r(d)$ is the rank of document $d$ in list $r$. Despite its simplicity, RRF was shown
to beat both the individual rankers and several more elaborate fusion schemes. That
balance of robustness and simplicity is why it was chosen to fuse the FAISS and BM25
results before the top-K articles are selected.

\section{Query Rewriting for Retrieval}

\hs Good retrieval depends not only on the index but on the query that reaches it. This
becomes a real problem when users write informally and the target documents are written
formally, which is exactly the situation in Cambodian legal search.
\citet{ma2023queryrewriting} examined this directly and showed that placing a rewriting
step before retrieval, so that the original question is reshaped into a form better suited
to the retriever, measurably improves the answers a RAG system produces. Their finding
supports one of the central design choices in this thesis, namely the dedicated query
rewriting stage that expands an everyday Khmer question into the formal legal vocabulary
likely to appear in the law itself.

\section{Legal Question Answering Systems}

\hs Applying language models to law brings its own risks, and recent work has tried to
measure them carefully. \citet{guha2023legalbench} assembled LegalBench, a broad benchmark
of legal reasoning tasks built collaboratively by legal and machine learning experts.
Their results make a point that is directly relevant here: general-purpose language models
can sound fluent on legal questions while still reasoning incorrectly. For a system meant
to give the public legal information, that gap between fluency and correctness is exactly
why this work grounds every answer in retrieved articles rather than trusting the model's
own knowledge. It also explains the inclusion of a grounding check that confirms the
retrieved articles are relevant before any answer is written.

\section{Agentic RAG and Reasoning Frameworks}

\hs A more recent line of work treats the system as something that decides what to do,
not just something that retrieves and generates in a fixed order. \citet{yao2023react}
introduced ReAct, which interleaves reasoning steps with concrete actions so that a
language model can fetch information and adjust course as it goes. The result was a clear
improvement over approaches that only reason or only retrieve. The pipeline in this thesis
shares that philosophy: it uses a directed state graph to judge a query's intent and the
quality of its retrieval before committing to an answer.

\hs The most directly related system is Self-RAG, proposed by \citet{asai2024selfrag}.
There, the model learns to decide when to retrieve, and then to critique both the
retrieved passages and its own draft answer using special reflection tokens. The grounding
check in this thesis pursues the same goal by simpler means, judging whether the retrieved
articles are relevant and, if they are not, rewriting the query and trying again before
allowing generation to proceed. Self-RAG shows that this kind of self-assessment is not
just a safeguard but a genuine source of quality, which is reassuring for a legal setting
where a wrong answer carries real consequences.

\section{Low-Resource Languages}

\hs Almost all of the work above was developed and tested in high-resource languages, and
Khmer is not among them. The SeaLLMs project by \citet{nguyen2023seallms} is one of the
few efforts to build large language models specifically for Southeast Asian languages,
Khmer included, and it documents the familiar difficulties of limited training data and
uneven tokenisation that such languages face. Their work helps explain why a purely
neural, embedding-based approach is risky for Khmer legal text. It also strengthens the
case made earlier for pairing dense retrieval with BM25, so that exact-term matching can
compensate in the many cases where the embeddings, trained on little Khmer, fall short.

\section{Summary of Related Work}

\hs Table~\ref{tab:literature_summary} brings these studies together and notes how each
one connects to the system proposed in this thesis.

\begin{table}[H]
\centering
\caption{Summary of Related Work}
\label{tab:literature_summary}
\begin{tabular}{p{4cm} p{3.5cm} p{6cm}}
\hline
\textbf{Study} & \textbf{Method} & \textbf{Relevance to This Work} \\
\hline
\citet{lewis2020rag} & RAG &
Foundation of the proposed RAG pipeline \\
\hline
\citet{karpukhin2020dpr} & Dense Retrieval (DPR) &
Basis for the FAISS dense retrieval component \\
\hline
\citet{robertson2009bm25} & BM25 &
Basis for the BM25 sparse retrieval component \\
\hline
\citet{johnson2021faiss} & FAISS &
Library used for dense vector indexing \\
\hline
\citet{cormack2009rrf} & RRF &
Fusion method for hybrid retrieval \\
\hline
\citet{ma2023queryrewriting} & Query Rewriting &
Motivates the query rewriting stage \\
\hline
\citet{guha2023legalbench} & LegalBench &
Motivates grounding over parametric knowledge \\
\hline
\citet{yao2023react} & ReAct &
Basis for the agentic pipeline design \\
\hline
\citet{asai2024selfrag} & Self-RAG &
Closest parallel to the grounding stage \\
\hline
\citet{nguyen2023seallms} & SeaLLMs &
Motivates hybrid retrieval for Khmer \\
\hline
\end{tabular}
\end{table}
